\chapter{Theoretical Framework}
\label{ch:framework}

\section{The mean-field model}

Take $N$ agents holding binary opinions $s_i = \pm 1$, coupled all-to-all with strength $J/N$,
in an external field $h$:
\begin{equation}
\mathcal{H} = -\frac{J}{N}\sum_{i<j} s_i s_j - h \sum_i s_i .
\end{equation}
The collective polarisation is $m = N^{-1}\sum_i s_i$, and in the thermodynamic limit it
satisfies the self-consistency condition
\begin{equation}
m = \tanh\!\left(\beta J m + h\right),
\label{eq:selfcons}
\end{equation}
with $\beta = 1/k_B T$ and $k_B = 1$ throughout. The field $h$ absorbs the factor $\beta$, so
it is dimensionless and directly comparable with $\beta J$.

The social reading of $\beta$ is an inverse noise level. Large $\beta$ describes a population
that copies its neighbours reliably; small $\beta$ one that decides close to at random.
Equation~\eqref{eq:selfcons} has a single root for $\beta J < 1$ at zero field and three for
$\beta J > 1$, of which the outer two are stable.

\section{The Landau expansion, with coefficients derived}

Write the free energy as a quartic in the order parameter,
\begin{equation}
F(m) = a\,m^2 + b\,m^4 - h\,m .
\label{eq:landau}
\end{equation}
The coefficients are not free. Requiring that $\partial F/\partial m = 0$ reproduce
\eqref{eq:selfcons} to cubic order fixes them. Inverting \eqref{eq:selfcons} gives
$\operatorname{artanh}(m) = \beta J m + h$, and expanding
$\operatorname{artanh}(m) = m + m^3/3 + O(m^5)$ yields
\begin{equation}
(1 - \beta J)\,m + \tfrac{1}{3} m^3 - h = 0 .
\end{equation}
Matching against $\partial F/\partial m = 2am + 4bm^3 - h$ gives
\begin{equation}
a = \frac{1 - \beta J}{2}, \qquad b = \frac{1}{12}.
\label{eq:coefficients}
\end{equation}

This derivation is stated in full because coefficient forms that do not follow from the
accompanying self-consistency condition appear in the applied literature, and the symptom is
subtle: the free-energy curve retains the correct qualitative double well, so the figure looks
right, while its minimum sits away from the $m^\ast$ the same work's root-finder reports. The
double well opens when $a < 0$, that is when $\beta J > 1$, recovering the transition already
implied by \eqref{eq:selfcons}.

The susceptibility follows as
\begin{equation}
\chi = \frac{\beta \operatorname{sech}^2(\beta J m^\ast + h)}
            {1 - \beta J \operatorname{sech}^2(\beta J m^\ast + h)},
\label{eq:chi}
\end{equation}
which diverges where the denominator vanishes and is undefined beyond it. It is reported only
where the denominator is positive.

\section{Numerical validation of the solver}

Every figure in this framework is drawn with one solver, so the solver is checked first
against results known exactly. Mean-field theory fixes three critical exponents:
$m^\ast \sim (\beta J - 1)^{\beta}$ with $\beta = 1/2$, $\chi \sim |1-\beta J|^{-\gamma}$ with
$\gamma = 1$, and $m \sim h^{1/\delta}$ on the critical isotherm with $\delta = 3$.

One methodological point matters. Solving \eqref{eq:selfcons} by direct iteration of
$m_{k+1} = \tanh(\beta J m_k + h)$ is adequate away from criticality but degrades exactly at
$\beta J = 1$, where the map's derivative reaches unity. The exponents are defined in that
limit, so fitting them to iterated values measures truncation error rather than physics: doing
so returns $\delta \approx 1.35$ against an exact $3$. The exponents are therefore fitted to
roots obtained by bracketing, and the iterative solver is checked against the bracketed one
away from criticality, where the largest absolute discrepancy over the sampled grid is
$1.15\times10^{-9}$.

Fitting within a window of $0.05$ in $|\beta J - 1|$ at 60 points recovers
\begin{equation}
\beta = 0.4872, \qquad \gamma = 1.0000, \qquad \delta = 3.0042,
\end{equation}
with absolute errors $0.0128$, $0.0000$ and $0.0042$. The residual in $\beta$ is a
finite-window effect: the power law is asymptotic and the quartic term contaminates the slope
as the window widens. The run fails if any exponent departs from its exact value by more than
a stated tolerance, so a solver regression stops the analysis rather than producing plausible
curves.

\begin{figure}[htbp]
\centering
\includegraphics[width=0.72\textwidth]{../../services/inference/data/paper1/F1_free_energy.pdf}
\caption{Free energy \eqref{eq:landau} at zero field with the derived coefficients
\eqref{eq:coefficients}, at $\beta J = 0.8$, $1.0$ and $1.4$. The second minimum appears only
above the critical coupling.}
\label{fig:free-energy}
\end{figure}

\begin{figure}[htbp]
\centering
\includegraphics[width=0.72\textwidth]{../../services/inference/data/paper1/F2_order_parameter.pdf}
\caption{Order parameter across the transition. The fitted exponent is $0.4872$ against an
exact $1/2$.}
\label{fig:order}
\end{figure}

\section{The spinodal field}

Above $\beta J = 1$ the free energy has two minima at zero field. Raising the field deepens
one and shallows the other; at a finite field the shallower minimum disappears. That value is
the spinodal field $h_c(\beta J)$, and it is the field a population in the metastable state
must experience before it is obliged to move.

It is located by bisection on the presence of a second stationary point, which keeps it
consistent with the solver used elsewhere rather than importing a closed form derived under
different approximations. It vanishes as $\beta J \to 1^{+}$ and grows monotonically, reaching
$0.53284$ at $\beta J = 2$.

\begin{figure}[htbp]
\centering
\includegraphics[width=0.72\textwidth]{../../services/inference/data/paper1/F5_phase_diagram.pdf}
\caption{Phase diagram in $(\beta J, h)$. The shaded region is where two opinion states
coexist; its upper boundary is the spinodal $h_c(\beta J)$.}
\label{fig:phase}
\end{figure}

\section{The bound}
\label{sec:bound}

Let an observable $X$ be proposed as the field through a linear map
\begin{equation}
h = \alpha X,
\end{equation}
with $\alpha$ an unknown coupling carrying the units. Any real observable is bounded: across
the content a population plausibly encounters, $X$ spans a finite range $\Delta X$, which is
an empirical property of the observable and the corpus rather than of the model. The largest
field the mechanism can deliver is therefore $\alpha\,\Delta X$, and requiring it to reach the
spinodal gives
\begin{equation}
\alpha \Delta X \;\ge\; h_c(\beta J)
\qquad \Longleftrightarrow \qquad
\alpha \;\ge\; \frac{h_c(\beta J)}{\Delta X}.
\label{eq:bound}
\end{equation}

Three properties of \eqref{eq:bound} carry the argument of this thesis.

It is \textbf{necessary and not sufficient}. Clearing it does not establish that content tips
populations; falling below it excludes the mechanism within the model's assumptions.

It \textbf{requires no fit}. $\Delta X$ is obtained by measuring the observable across a
corpus, which is far cheaper than estimating a coupling to collective behaviour, and $h_c$
comes from the model.

It \textbf{converts a null into a constraint}. A study that estimates $\alpha$ and obtains an
interval containing zero has, alone, only the statement that no coupling was detected.
Combined with \eqref{eq:bound}, the same interval says the coupling would have to exceed
$h_c(\beta J)/\Delta X$ for the mechanism to operate, and that the measurement excludes that
range.

\begin{figure}[htbp]
\centering
\includegraphics[width=0.72\textwidth]{../../services/inference/data/paper1/F6_alpha_required.pdf}
\caption{Coupling required before a content observable of spread $\Delta X$ can drive the
population across the spinodal.}
\label{fig:alpha-required}
\end{figure}

\section{A screening criterion}

The bound supports a procedure that runs before data collection.

\begin{enumerate}
\item State the observable and the map $h = \alpha X$ explicitly, with units.
\item Measure or bound $\Delta X$ on a modest sample. This needs no outcome variable.
\item Choose the range of $\beta J$ the population plausibly occupies and read $h_c(\beta J)$.
\item Evaluate $\alpha_{\min} = h_c(\beta J)/\Delta X$ and ask whether a coupling of that size
is defensible on independent grounds.
\item If it is not, the mechanism is excluded within the model and the expensive measurement
is unnecessary. If it is, proceed, and report $\Delta X$ alongside the estimate of $\alpha$ so
the bound can be applied to the result.
\end{enumerate}

Step 4 is the weakest link, since $\alpha$ has no independent scale in most applications. The
criterion remains informative in relative terms: two candidate observables can be ranked by
$\alpha_{\min}$, and an observable whose spread demands an implausibly large coupling is a poor
choice of field regardless of how well it correlates with anything.

\section{Assumptions}

The framework assumes all-to-all coupling, a single homogeneous population, a static field and
binary opinions. Network structure and heterogeneous susceptibility generally lower $h_c$, so
\eqref{eq:bound} is a bound within this model rather than a universal one. The linear map
$h = \alpha X$ is the simplest possible coupling and is itself an assumption; a saturating map
would change the relation between $\Delta X$ and the maximum attainable field. Finally, the
spinodal is the threshold relevant to a population trapped in a metastable state; a population
at coexistence can be moved by a smaller field given enough time, so the bound demands more of
$\alpha$ than a nucleation argument would.
